\documentclass[11pt,a4paper]{article}
\usepackage[T1]{fontenc}
\usepackage[utf8]{inputenc}
\usepackage{lmodern}
\usepackage{amsmath,amssymb,amsthm}
\usepackage[margin=1.15in]{geometry}
\usepackage[hidelinks]{hyperref}
\emergencystretch=3em

\newtheorem{theorem}{Theorem}
\renewcommand{\thetheorem}{\Alph{theorem}}
\newtheorem{corollary}[theorem]{Corollary}
\newtheorem{lemma}{Lemma}[section]
\theoremstyle{definition}
\newtheorem{remark}[lemma]{Remark}

\newcommand{\A}{\mathbb{A}}
\newcommand{\PP}{\mathbb{P}}
\newcommand{\Z}{\mathbb{Z}}
\newcommand{\cO}{\mathcal{O}}
\newcommand{\Spec}{\operatorname{Spec}}
\newcommand{\Proj}{\operatorname{Proj}}
\newcommand{\Frac}{\operatorname{Frac}}
\newcommand{\Gal}{\operatorname{Gal}}
\newcommand{\divi}{\operatorname{div}}

\title{A wild quasi-\'etale double cover over a single smooth genus-two
exceptional curve in characteristic two\thanks{AI-produced. The example and
its proof are due to a run of GPT-6 Astra (OpenAI, ChatGPT Pro mode);
Corollary~\ref{cor:C}\,(ii)--(iii), the proof of Lemma~\ref{lem:ext} given
here and the text are due to Claude Opus~5.5 (Anthropic) at effort max, and
Lemma~\ref{lem:rational} to a second agent of the same model and effort; see
Section~\ref{sec:provenance}. No human has checked this text. To the extent
that any rights exist in it, they are dedicated to the public domain under
CC0~1.0.}}
\author{}
\date{25 November 2026}

\begin{document}
\maketitle

\begin{abstract}
Over an algebraically closed field of characteristic two, the surface
$b^2+(t^5+t^2a^2)b+a^5=0$ has one singular point, and its minimal resolution
has one exceptional curve $C$: the smooth curve of genus two with affine
equation $y^2+y=x^5$, with $C^2=-1$. We show that this surface has a
connected Galois double cover, \'etale away from the singular point, whose
ramification along $C$ is fierce: ramification index one and a purely
inseparable residue extension of degree two. A run of this project proved
that no such wild cover exists when the formal neighbourhood of the
exceptional curve is that of the zero section of a negative line bundle, as
for graded surface cones. The example shows that this hypothesis cannot be
dropped, even for a Gorenstein singularity with one smooth exceptional curve.
The weighted homogeneous surface $b^2+t^5b+a^5=0$ has the same exceptional
curve, with the same self-intersection, and every connected Galois cover of
it that is \'etale away from the singular point is tame along that curve; so
the two germs are not formally isomorphic. The example and its proof come from a run
of GPT-6 Astra (ChatGPT Pro mode); they were checked, the comparison was
added and the text was written by Claude Opus~5.5 at effort max.
\end{abstract}

\section{Results}\label{sec:results}

Throughout, $k$ is an algebraically closed field. For a normal surface germ,
a cover \emph{\'etale on the punctured germ} is a finite cover by a normal
germ that is \'etale away from the closed point. Let $R=\bigoplus_{d\ge0}R_d$
be a graded ring with $R_0=k$ and let $A$ be its completion at
$R_+=\bigoplus_{d>0}R_d$; the \emph{radial valuation} of $\Frac(A)$ is
$v(f)=\min\{d: f_d\ne0\}$ for $0\ne f=\sum_d f_d\in A$.

\begin{theorem}[{\cite[Lemma~4.16 and Proposition~4.17, pp.~105--107]{SA12}}]
\label{thm:A}
Let $k$ have characteristic $p>0$, let $R=\bigoplus_{d\geq0}R_d$ be a
finitely generated normal graded $k$-domain of dimension two with $R_0=k$,
and let $A$ be its completion at $R_+$. Let $B/A$ be the normal extension
belonging to a connected finite Galois cover of the punctured spectrum of
$A$. Assume the positive grading is effective, so that its occurring degrees
have greatest common divisor one. Then every extension to $\Frac(B)$ of the
radial divisorial valuation of $\Frac(A)$ is tame. Thus a normal surface germ
whose completed resolution along its exceptional curve is the completion of
the zero section in a negative line bundle over a smooth proper curve has no
wild inertia at that exceptional valuation in any finite cover \'etale on the
punctured germ.
\end{theorem}

From now on $k$ has characteristic two. Put
\begin{align*}
F&=b^2+(t^5+t^2a^2)\,b+a^5, & F_0&=b^2+t^5b+a^5,\\
H&=d^2+(c^4+t^3+t^4c^2)\,d+tc^5+t^2c^7+t^4c^6+t^6c^5,
\end{align*}
let $S$, $S_0$ and $T$ be the affine surfaces $F=0$, $F_0=0$ in
$\A^3_{a,b,t}$ and $H=0$ in $\A^3_{c,d,t}$, and let $o$ be the origin. Let
$\varphi\colon T\to S$ and $\sigma\colon T\to T$ be given by
\begin{align*}
\varphi^*a&=c^2+t^2c, & \varphi^*b&=t^2d+c^5, & \varphi^*t&=t,\\
\sigma^*c&=c+t^2, & \sigma^*d&=d+c^4+t^6c+t^8, & \sigma^*t&=t.
\end{align*}

\begin{theorem}[{\cite[Example~4.60, p.~143]{SA12}}]\label{thm:B}
Let $k$ have characteristic two.
\begin{enumerate}
\item[(i)] $S$ is a normal Gorenstein surface whose only singular point is
$o$. The weighted blow-up of $\A^3$ with weights $(2,5,1)$ on $(a,b,t)$
restricts to a resolution $f\colon Y\to S$ which is an isomorphism over
$S\setminus\{o\}$. Its exceptional fibre is one smooth curve $C$, the smooth
projective model of $y^2+y=x^5$, of genus two, with $C^2=-1$. The
discrepancy of $C$ is $-3$.
\item[(ii)] $T$ is a normal surface, $\varphi$ is finite of degree two and
\'etale over $S\setminus\{o\}$, $\varphi^{-1}(o)$ is one point $o'$, and
$\sigma$ generates $\Gal(k(T)/k(S))\cong\Z/2$. The surface $T$ is singular
at $o'$, and $o'$ is not a rational singularity; $T$ is smooth elsewhere.
\item[(iii)] The valuation $v_C$ of $C$ has a unique extension $w$ to $k(T)$.
Its ramification index is one, its residue extension is purely inseparable of
degree two, its different exponent is one, and $\sigma$ acts trivially on its
residue field. So the inertia group at $w$ is all of $\Z/2$: the cover
$\varphi$ is wildly ramified along $C$.
\end{enumerate}
In the coordinates $x=a/t^2$ on $Y$ and $z=c/t$ on $T$ the cover is
$z^2+tz=x$.
\end{theorem}

\begin{corollary}\label{cor:C}
In the notation above:
\begin{enumerate}
\item[(i)] The completion of $Y$ along $C$ is not isomorphic to the
completion of the total space of a negative line bundle on a smooth proper
curve along its zero section.
\item[(ii)] $S_0$ is a normal surface whose only singular point is $o$; the
same weighted blow-up gives a resolution of $S_0$ with one exceptional curve
$C_0\cong C$, with $C_0^2=-1$; and every connected finite Galois cover of the
punctured spectrum of $\widehat\cO_{S_0,o}$ is tame along $C_0$.
\item[(iii)] The complete local rings $\widehat\cO_{S,o}$ and
$\widehat\cO_{S_0,o}$ are not isomorphic.
\end{enumerate}
\end{corollary}

\paragraph{Why it matters.} Starting from Grothendieck's notes on the
fundamental group in cote~29 of his Montpellier archive, the run
\cite{SA12} studied the topological finite generation of
$\pi_1(X\setminus Y)$, and of its quotient by the wild inertia at the
codimension-one components of $Y$, for $X$ normal and proper and $Y\subset
X$ closed \cite[\S4.1]{SA12}. Where $Y$ has isolated points at which $X$ is
singular, one resolves $X$ there (at a regular point, a cover \'etale off the
point extends, by purity), and a cover \'etale over $X\setminus Y$ may then
be wildly ramified along the new exceptional curves, which a tame fundamental group of the
resolution does not see. Theorem~\ref{thm:A} excludes this when the
exceptional curve has a linear formal neighbourhood. Theorem~\ref{thm:B}
shows that the hypothesis is needed: one smooth exceptional curve of genus
two, on a Gorenstein surface, already carries a fierce cover. Corollary~\ref{cor:C} shows that the dual graph and the
self-intersection do not decide the matter: $S$ and $S_0$ have the same
exceptional curve $C$ with $C^2=-1$, and only $S$ has a wild cover.

\paragraph{Credit.} Theorem~\ref{thm:A} and Theorem~\ref{thm:B}, with their
proofs, are due to a run of GPT-6 Astra (OpenAI, ChatGPT Pro mode) of this
project \cite{SA12}. Corollary~\ref{cor:C}\,(i) makes precise the run's remark
that the hypothesis of Theorem~\ref{thm:A} cannot be discarded.
Corollary~\ref{cor:C}\,(ii)--(iii) and the proof of Lemma~\ref{lem:ext}
given here are due to the writer, Claude Opus~5.5 at effort max;
Lemma~\ref{lem:rational} is due to the note's independent checker, Claude
Opus~5.5 at effort max. Details are in Section~\ref{sec:provenance}.

\paragraph{Context: what is known and what is new.} Known: wild covers of
surface singularities, \'etale off the closed point, exist in positive
characteristic: Artin's wild $\Z/2$-actions on $k[[u,v]]$ \cite{Ar75} and
coverings of rational double points \cite{Ar77}, the wild $\Z/p$-quotients
of $k[[u,v]]$ of Lorenzini--Schr\"oer \cite{LS20,LS23}, the wild quotients
of Ito--Schr\"oer \cite{IS} by special $p$-groups acting diagonally on
products of Hermitian curves, and Matsumoto's $\Z/p$-coverings of rational
double points by regular germs \cite[\S6.2 and Lemma~6.6]{Ma}. In all of them
the source of the cover is regular or its base is a rational double point,
and then every exceptional curve over the base is rational
(Lemma~\ref{lem:rational}; for quotients of regular germs see also
\cite[p.~5]{IS}). Here $C$ has genus two, and $T$ is not a rational
singularity. Cutkosky and Srinivasan \cite[Thm.~1(b)]{CS} show that if the
prime-to-$p$ completion of the local fundamental group is finite, the
exceptional curves of a good resolution are rational and form a tree; here
$C$ has genus two, so that completion is infinite. For strongly $F$-regular
singularities the local fundamental group is finite of order prime to $p$
\cite[Thm.~A]{CST}, so their covers are tame; $S$ lies far outside that
class, since its discrepancy $-3$ is below $-1$. New, to our knowledge after
the search described in Section~\ref{sec:provenance}: the example of
Theorem~\ref{thm:B} and Corollary~\ref{cor:C}.

\paragraph{Scope.} We treat one example in characteristic two. We do not
claim that every non-linear formal neighbourhood carries a wild cover. The
point $o'$ of $T$ is a singular, non-rational double point
(Lemma~\ref{lem:rational}). The run also resolves $T$ (one curve, the
Frobenius twist of $C$, with self-intersection $-2$) and computes the maximal pro-$2$ quotients of the
local fundamental groups of $T$ and $S$ \cite[Prop.~4.61, Cor.~4.62]{SA12};
they are not used here. The example bears on tameness only, not on the
finite generation of any fundamental group.

\paragraph{Proof overview.} $S$ is a hypersurface with an isolated
singularity. The weighted blow-up has two smooth charts covering the
exceptional fibre, $y^2+(1+tx^2)y+x^5=0$ and
$\xi^2(1+\xi+\rho\tau^2)+\tau^5=0$; the first shows $C$ and its genus, and
the divisor of $t$ on the second gives $C^2=-1$. The function $x=a/t^2$ is a
unit along $C$ whose residue is not a square in $k(C)$, so $z^2+tz=x$ defines
a separable quadratic extension in which $t$ stays a uniformizer and the
residue field becomes $k(C)(\sqrt{x})$. The substitution $c=tz$,
$d=(b+c^5)/t^2$ turns $F$ into $t^4H$, which presents the normalization $T$;
the cover is \'etale where $t\ne0$ by the derivative $t^2$, and where $t=0$
by a Jacobian computation. Corollary~\ref{cor:C} follows from
Theorem~\ref{thm:A}, applied to $S$ and to the graded ring of $S_0$.

\section{Proofs}\label{sec:proofs}

\subsection{The surfaces and their resolution}

\begin{lemma}\label{lem:normal}
$F$ and $F_0$ are irreducible, and $S$ and $S_0$ are normal Gorenstein
surfaces whose only singular point is $o$.
\end{lemma}

\begin{proof}
Both polynomials are monic of degree two in $b$ over the factorial ring
$k[a,t]$; a factorization would give a root in $k(t)(a)$. Put
$h=t^5+t^2a^2$ (resp.\ $h=t^5$) and $w=b/h$; a root gives
$a^5/h^2=w^2+w$ with $w\in k(t)(a)$. For the valuation $-\deg_a$ at $a=\infty$
over $k(t)$, the left side has value $-5+2\deg_a h$, which is $-1$
(resp.\ $-5$), an odd negative number, while $w^2+w$ has value $2v(w)$ if
$v(w)<0$ and a value $\ge0$ otherwise. So $F$ and $F_0$ are irreducible. In
characteristic two,
\[
F_a=a^4,\quad F_b=t^2(t^3+a^2),\quad F_t=t^4b;\qquad
(F_0)_a=a^4,\quad (F_0)_b=t^5,\quad (F_0)_t=t^4b.
\]
On either surface, $a=0$ forces $t=0$ through the second derivative and then
$b^2=0$; so the singular locus is $\{o\}$. A hypersurface is Gorenstein and
satisfies Serre's condition $(S_2)$, and $(R_1)$ holds; so both are normal.
\end{proof}

\begin{lemma}\label{lem:resolution}
The weighted blow-up of Theorem~\ref{thm:B}\,(i) restricts to resolutions
$f\colon Y\to S$ and $f_0\colon Y_0\to S_0$ which are isomorphisms over the
complement of $o$. The exceptional fibre of each is a single smooth curve
isomorphic to the smooth projective model $C$ of $y^2+y=x^5$, of genus two
and self-intersection $-1$. The discrepancy of $C$ in $Y$ is $-3$.
\end{lemma}

\begin{proof}
Let $I_m\subset k[a,b,t]$ be spanned by the monomials $a^ib^jt^l$ of weight
$2i+5j+l\ge m$, and $X=\Proj\bigoplus_{m\ge0}I_ms^m$. The Rees algebra is
generated by $as,as^2,bs,\dots,bs^5,ts$, so $X\to\A^3$ is projective; it is
an isomorphism over $\A^3\setminus\{o\}$, since there one coordinate is a unit
and its powers lie in every $I_m$. Its exceptional divisor is
$E=\Proj\bigoplus_m I_m/I_{m+1}=\PP(2,5,1)$, and $X$ is covered by
$D_+(ts)$, $D_+(as^2)$ and $D_+(bs^5)$, since $D_+(as)\subset D_+(as^2)$ and
$D_+(bs^j)\subset D_+(bs^5)$.

\emph{The chart $D_+(ts)$} is $\Spec k[x,y,t]$ with $x=a/t^2$, $y=b/t^5$, and
$E=\{t=0\}$ on it. Since $F(t^2x,t^5y,t)=t^{10}G_1$ with
$G_1=y^2+(1+tx^2)y+x^5$, and $F_0(t^2x,t^5y,t)=t^{10}(y^2+y+x^5)$, the strict
transforms there are $G_1=0$ and $y^2+y+x^5=0$. They are smooth: $\partial_y
G_1=1+tx^2$ and $\partial_xG_1=x^4$ have no common zero, and likewise for the
second. They meet $E$ in $C^{\rm aff}=\{y^2+y+x^5=0\}\subset\A^2$.

\emph{The chart $D_+(as^2)\cap D_+(bs^5)$.} Put $\rho=b/a^2$, $\xi=a^5/b^2$,
$\tau=ta^2/b$, so that $a=\rho^2\xi$, $b=\rho^5\xi^2$, $t=\rho\tau$. The
monomial $\rho^\alpha\xi^\beta\tau^\gamma$ equals
$a^{-2\alpha+5\beta+2\gamma}b^{\alpha-2\beta-\gamma}t^\gamma$, of weight
$\alpha$; the exponent matrix has determinant $-1$. The ring of this chart is
spanned by the Laurent monomials in $a,b$ times powers of $t$ of non-negative
weight, so it is $k[\rho,\tau,\xi,\xi^{-1}]$, with $E=\{\rho=0\}$. Since
$F(\rho^2\xi,\rho^5\xi^2,\rho\tau)=\rho^{10}\xi^2G_2$ with
$G_2=\xi^2(1+\xi+\rho\tau^2)+\tau^5$, and the same with
$\xi^2(1+\xi)+\tau^5$ for $F_0$, the strict transforms are $G_2=0$ and
$\xi^2(1+\xi)+\tau^5=0$. They are smooth, as $\partial_\xi$ of either
equation is $\xi^2$, a unit.

\emph{The exceptional fibre.} The strict transform $Y$ lies in the zero locus
of $Fs^{10}$, whose intersection with $E$ is the zero locus of the initial
form $b^2+t^5b+a^5$ in $\PP(2,5,1)$; the same holds for $Y_0$. At a point of
that locus with $t=0$ we have $b^2=a^5$, so $a\ne0\ne b$, and the point lies
in the second chart, with $\rho=\tau=0$ and $\xi^2(1+\xi)=0$, hence $\xi=1$.
So the exceptional fibre of $Y$, and that of $Y_0$, is
$C^{\rm aff}\cup\{P_\infty\}$ with $P_\infty=(\rho,\tau,\xi)=(0,0,1)$; it is a
proper curve, smooth by the two charts (at $P_\infty$ the curve is
$\xi^2(1+\xi)+\tau^5=0$ with $\partial_\xi=1$), and irreducible since
$y^2+y+x^5$ is (the argument of Lemma~\ref{lem:normal} with $h=1$). Hence it
is the smooth projective model $C$ of $y^2+y=x^5$, and $Y$, $Y_0$ are smooth
along it.

\emph{Genus.} Let $\pi\colon C\to\PP^1_x$ be the double cover. Over $x\ne\infty$,
$\pi_*\cO_C$ is free on $1,y$. At infinity put $u=1/x$ and $v=y/x^3$: then
$v^2+u^3v+u=0$, which is smooth at $(0,0)$ since its $u$-derivative is
$1+u^2v$, and $\pi_*\cO_C$ is free on $1,v$. As $y=u^{-3}v$, we get
$\pi_*\cO_C\cong\cO\oplus\cO(-3)$, so $\chi(\cO_C)=1+(1-3)=-1$ and $g(C)=2$
(this is also the Artin--Schreier formula \cite[Prop.~3.7.8]{Sti}).

\emph{Self-intersection.} On $D_+(ts)$ the function $t$ is a local equation
of $C$. On the second chart $t=\rho\tau$, where $\rho=0$ is $C$ and $\tau=0$
is a curve $D$ meeting $C$ only at $P_\infty$; there $(\rho,\tau)$ are local
parameters, because $\partial_\xi G_2$ is a unit. So
$\divi_Y(t)=C+D$ with $C\cdot D=1$, and $0=C\cdot\divi_Y(t)=C^2+1$. The same
holds on $Y_0$.

\emph{Discrepancy.} $S$ is Gorenstein, so $K_Y=f^*K_S+\alpha C$ with
$\alpha\in\Z$. Adjunction gives $K_Y\cdot C=2g(C)-2-C^2=3$, hence
$\alpha C^2=3$ and $\alpha=-3$.
\end{proof}

\subsection{The double cover}

Let $K=k(S)$, $x=a/t^2\in K$, and let $\cO=\cO_{Y,C}$ be the local ring of
$Y$ at the generic point of $C$: a discrete valuation ring with uniformizer
$t$ and residue field $k(C)=k(x,y)$, $y^2+y=x^5$, in which $x$ is a unit.

\begin{lemma}\label{lem:ext}
Let $L=K[z]/(z^2+tz+x)$. Then $L$ is a separable quadratic field extension
of $K$, the integral closure of $\cO$ in $L$ is $\cO'=\cO[z]$, a discrete
valuation ring with uniformizer $t$ and residue field $k(C)(\sqrt{x})$,
purely inseparable of degree two over $k(C)$, and the different of
$\cO'/\cO$ is $(t)$.
\end{lemma}

\begin{proof}
The residue $\bar x$ of $x$ is not a square in $k(C)$: the extension
$k(C)/k(\bar x)$ is separable, since $y^2+y+\bar x^5$ has derivative $1$
in $y$, so $d\bar x\ne0$ in $\Omega_{k(C)/k}$, whereas squares have zero
differential. Let $g=z^2+tz+x$ and $\cO'=\cO[z]/(g)$, free of rank two over
$\cO$. Then $\cO'/t\cO'=k(C)[z]/(z^2+\bar x)=k(C)(\sqrt{\bar x})$ is a field.
So $\cO'$ is local with principal maximal ideal $t\cO'$, and $t$ is a
non-zerodivisor; $\cO'$ is a one-dimensional regular local ring, a discrete
valuation ring. In particular $\cO'$ is a domain, so $g$ is irreducible over
$K$ and $L=\Frac(\cO')$ is a field; $g'(z)=t\ne0$, so $L/K$ is separable.
Being normal and finite over $\cO$, $\cO'$ is the integral closure of $\cO$
in $L$. Since $\cO'=\cO[z]$ is monogenic, its different is
$(g'(z))=(t)$ \cite[III, \S6, Cor.~2]{Se}.
\end{proof}

\begin{lemma}\label{lem:T}
$T$ is a normal surface with function field $L$, $\varphi$ is finite of
degree two and identifies $T$ with the normalization of $S$ in $L$ through
$c=tz$, $d=(b+c^5)/t^2$, and $\sigma$ generates $\Gal(L/K)$, with
$\sigma(z)=z+t$.
\end{lemma}

\begin{proof}
Direct expansion in characteristic two gives the identity
\begin{equation}\label{eq:I}
F(c^2+t^2c,\;t^2d+c^5,\;t)=t^4H(c,d,t)
\end{equation}
in $k[c,d,t]$; it is also checked by machine (Section~\ref{sec:provenance}).
So $\varphi$ is well defined. Let $A_S=k[a,b,t]/(F)$ and $A_T=k[c,d,t]/(H)$.
Over $A_S$, $c$ satisfies $c^2+t^2c=a$ and $d$ satisfies the monic equation
$H=0$ over $k[c,t]$, so $A_T$ is finite over $A_S$.

Inverting $t$, the maps $a\mapsto c^2+t^2c$, $b\mapsto t^2d+c^5$ and
$d\mapsto(b+c^5)/t^2$ are inverse isomorphisms
$A_S[t^{-1}][c]/(c^2+t^2c+a)\cong A_T[t^{-1}]$, by \eqref{eq:I}. The left side
is free of rank two over the domain $A_S[t^{-1}]$, and its generic fibre is
$K[c]/(c^2+t^2c+a)\cong L$ via $c=tz$, a field by Lemma~\ref{lem:ext}; so
$A_T[t^{-1}]\subset L$ is a domain. As $A_T$ is free over $k[c,t]$ on $1,d$,
$t$ is a non-zerodivisor on it, $A_T\subset A_T[t^{-1}]$, and $A_T$ is a
domain with fraction field $L$.

$T$ is smooth where $t\ne0$: there it is \'etale over $S\setminus\{t=0\}$,
which is smooth, because $c^2+t^2c+a$ has $c$-derivative $t^2$. Where $t=0$,
$H=d(d+c^4)$ and $\partial_dH=c^4$; this vanishes only for $c=0$, and then
$H=0$ forces $d=0$. So the only possible singular point is the origin, and
the hypersurface $T$ is normal by Serre's criterion. Being normal, finite
over $S$ and with fraction field $L$, $T$ is the normalization of $S$ in $L$.

The map $\sigma$ preserves $\varphi^*a$, $\varphi^*b$ and $t$: in
characteristic two $(c+t^2)^2+t^2(c+t^2)=c^2+t^2c$ and
$t^2(d+c^4+t^6c+t^8)+(c+t^2)^5=t^2d+c^5$. By \eqref{eq:I},
$t^4H\circ\sigma=t^4H$, so $\sigma$ preserves $H$; applied twice it is the
identity, and it is not the identity. So $\sigma$ is an automorphism of $T$
over $S$ of order two, which generates $\Gal(L/K)\cong\Z/2$. From $c=tz$ and
$\sigma(c)=c+t^2$ we get $\sigma(z)=z+t$.
\end{proof}

\begin{lemma}\label{lem:etale}
$\varphi^{-1}(o)$ is the origin of $T$, and $\varphi$ is \'etale over
$S\setminus\{o\}$.
\end{lemma}

\begin{proof}
Over $o$ we have $c^2=0$ and $t=0$, so $c=0$ and $H=d^2=0$. Over $t\ne0$,
$\varphi$ is \'etale by the proof of Lemma~\ref{lem:T}. Let $P\in
S\setminus\{o\}$ have $t=0$. Then $b^2=a^5$ forces $a\ne0$; since $F_a=a^4$
is a unit at $P$, $(b,t)$ are \'etale coordinates of $S$ at $P$. At a point
$Q$ over $P$ we have $c^2=a\ne0$, and $\partial_dH=c^4$ is a unit at $Q$, so
$(c,t)$ are \'etale coordinates of $T$ at $Q$. From $b=t^2d+c^5$ we get
$db\wedge dt=c^4\,dc\wedge dt$ at $Q$, where $t=0$, and $c^4\ne0$. So the
Jacobian of $\varphi$ is a unit at $Q$, and $\varphi$ is \'etale there.
\end{proof}

\begin{lemma}\label{lem:rational}
Let $T'\to S'$ be a finite surjective morphism of normal two-dimensional
excellent local schemes with algebraically closed residue fields. If $T'$ is
a rational singularity, for instance regular, then every exceptional curve of
every resolution of $S'$ is rational. In particular the point $o'$ of $T$
over $o$ is a singular, non-rational double point.
\end{lemma}

\begin{proof}
Let $Y'\to S'$ be a resolution and $E$ an exceptional curve. The
normalization $Y'_T$ of $Y'$ in $k(T')$ is finite over $Y'$ and maps to
$T'$, the normalization of $S'$ in $k(T')$; so it is proper and birational
over $T'$. As $Y'_T\to Y'$ is finite and surjective, some irreducible curve
$E_T\subset Y'_T$ maps onto $E$, and it lies over the closed point of $T'$.
Let $Z\to Y'_T$ be a resolution and $E_Z$ the strict transform of $E_T$.
Then $Z\to T'$ is a resolution, so $H^1(Z,\cO_Z)=0$ as $T'$ is rational. The
map $H^1(Z,\cO_Z)\to H^1(E_Z,\cO_{E_Z})$ is onto, since coherent cohomology
of $Z$ vanishes in degree two, the fibres of $Z\to T'$ having dimension at
most one. So the integral proper curve $E_Z$ has arithmetic genus zero,
$E_Z\cong\PP^1$, and $E$, which is dominated by $E_Z$, is rational by
L\"uroth's theorem. Apply this to $\Spec\cO_{T,o'}\to\Spec\cO_{S,o}$, which is
finite since $\varphi^{-1}(o)=\{o'\}$, and to the curve $C$ of genus two:
$o'$ is not a rational singularity. Directly, every monomial of $H$ has
degree at least two, and the quadratic part of $H$ is $d^2$; so $o'$ is a
singular double point of $T$.
\end{proof}

\begin{proof}[Proof of Theorem~\ref{thm:B}]
Part (i) is Lemmas~\ref{lem:normal} and~\ref{lem:resolution}, and part (ii)
is Lemmas~\ref{lem:T}, \ref{lem:etale} and~\ref{lem:rational}. For (iii), Lemma~\ref{lem:ext}
shows that $v_C$ has the unique extension $w$ given by $\cO'$, with
ramification index one, the purely inseparable residue extension
$k(C)(\sqrt{x})/k(C)$ of degree two and different exponent one. Since
$\sigma(z)=z+t\equiv z\pmod{t\cO'}$, $\sigma$ acts trivially on the residue
field. So the inertia group of $w$ is $\Gal(L/K)\cong\Z/2$, a $2$-group in
characteristic two: the ramification along $C$ is wild.
\end{proof}

\subsection{Proof of Corollary~\ref{cor:C}}

\emph{The cover on the completed germ.} Let $\widehat S=\Spec\widehat\cO_{S,o}$
and $\widehat Y=Y\times_S\widehat S$. As $\varphi^{-1}(o)$ is one point $o'$,
$T\times_S\widehat S=\Spec\widehat\cO_{T,o'}$ with
$\widehat\cO_{T,o'}=k[[c,d,t]]/(H)$, a two-dimensional complete intersection,
whose punctured spectrum is connected by Hartshorne's connectedness theorem
(depth at least two) \cite[\S2]{Ha}. So $\varphi$ gives a connected Galois cover of the
punctured spectrum of $\widehat\cO_{S,o}$ with group $\Z/2$. The morphism
$\widehat Y\to Y$ is flat, and its fibre over the generic point of $C$ is that
point, since $C$ lies over the closed point and
$\widehat\cO_{S,o}\otimes_{\cO_{S,o}}k=k$. So $\cO_{Y,C}\to\cO_{\widehat Y,C}$
is a flat local map of discrete valuation rings with the same uniformizer $t$
and the same residue field $k(C)$, and the proof of Lemma~\ref{lem:ext}
applies verbatim over $\cO_{\widehat Y,C}$: the completed cover is wild along
$C$. Moreover $\widehat Y$ is regular (\cite[Thm.~23.7]{Mat}: $Y$ is regular
and $\widehat Y\to Y$ is a base change of the regular morphism
$\widehat S\to\Spec\cO_{S,o}$, $S$ being excellent), proper and birational
over $\widehat S$, with exceptional curve $C$. It is the minimal resolution
of $\widehat S$: its factorization through the minimal resolution
\cite[p.~155, (A)]{Li78} cannot contract $C$, as $\widehat S$ is singular, so
it is finite and birational onto a normal surface, hence an isomorphism. In
the same way $\widehat Y_0=Y_0\times_{S_0}\Spec\widehat\cO_{S_0,o}$ is the
minimal resolution of $\Spec\widehat\cO_{S_0,o}$.

(i) If the completion of $Y$ along $C$ were that of the zero section of a
negative line bundle over a smooth proper curve, Theorem~\ref{thm:A} would
exclude wild inertia at $v_C$ in every finite cover \'etale on the punctured
germ, while $\varphi$ is such a cover with wild inertia at $v_C$.

(ii) By Lemmas~\ref{lem:normal} and~\ref{lem:resolution}, $S_0$ is normal
with only singular point $o$, and $Y_0$ is a resolution with one exceptional
curve $C_0\cong C$ of self-intersection $-1$. The ring
$R=k[a,b,t]/(F_0)$, graded by the weights $(2,5,1)$, is a finitely generated
normal domain of dimension two with $R_0=k$, effectively graded since $t$ has
degree one, and $\widehat\cO_{S_0,o}$ is its completion at $R_+$. Its radial
valuation is the valuation $\hat v_{C_0}$ of $C_0$ on $\widehat Y_0$, which
is regular as $\widehat Y$ is; $\hat v_{C_0}$ restricts to $v_{C_0}$ on
$k(S_0)$. Indeed, a homogeneous $f\ne0$ of degree $e$ equals $t^ef(x,y,1)$ on
the chart $D_+(ts)$, and $f(x,y,1)$ does not vanish on $C_0$: otherwise
$y^2+y+x^5$ would divide it, the quotient would have weighted degree at most
$e-10$ for the weights $(2,5)$ on $(x,y)$, and so $F_0$ would divide $f$. So
$v_{C_0}(f)=e$. For $f=\sum_df_d$ in the completion, the part of degree at
least $N$ lies in the ideal generated by the homogeneous elements of degree
at least $N$, so its value under $\hat v_{C_0}$ is at least $N$; hence
$\hat v_{C_0}(f)=\min\{d:f_d\ne0\}$. Theorem~\ref{thm:A} gives the tameness.

(iii) An isomorphism $\widehat\cO_{S,o}\cong\widehat\cO_{S_0,o}$ lifts to an
isomorphism of the minimal resolutions $\widehat Y$ and $\widehat Y_0$, which
are unique up to unique isomorphism \cite[p.~155, (A)]{Li78}, and it carries
$\hat v_C$ to $\hat v_{C_0}$, the valuations of their only exceptional
curves. It would carry the wild cover of the punctured
spectrum of $\widehat\cO_{S,o}$ to one of $\widehat\cO_{S_0,o}$ that is wild
along $C_0$, contradicting (ii). \qed

\section{Provenance and checks}\label{sec:provenance}

This paper is AI-produced. No human has proved, written or checked any part
of it.

\paragraph{What is whose.} Theorems~\ref{thm:A} and~\ref{thm:B} and their
proofs are due to the run RESEARCH\_SA12 of GPT-6 Astra (OpenAI, ChatGPT Pro
mode) of this project, which ended on 20 November 2026 \cite{SA12}. The
proof of Theorem~\ref{thm:B} above follows the run's, except for
Lemma~\ref{lem:ext}, where the run shows that the extension is non-trivial by
a \v{C}ech torsor on a family of curves. Corollary~\ref{cor:C}\,(ii)--(iii),
that proof of Lemma~\ref{lem:ext} and the text are due to the writer, Claude
Opus~5.5 at effort max (agent a3c8d7de, 25 November 2026).
Lemma~\ref{lem:rational}, on which the last sentence of
Theorem~\ref{thm:B}\,(ii) rests, and the argument identifying the radial
valuation in Corollary~\ref{cor:C}\,(ii) are due to the independent checker,
Claude Opus~5.5 at effort max (agent a65dbdc2, 25 November 2026).

\paragraph{Checks.} The run's own audits (GPT-6 Astra, ChatGPT Pro mode)
passed. At the intake of the run (21--22 November 2026) the writer made a
targeted check of Theorem~\ref{thm:A} and re-proved Theorem~\ref{thm:B} by
hand, including \eqref{eq:I} term by term; it reran the run's certificate
(26 exact assertions over $\mathbb{F}_2$) on 22 and 25 November. On 24 November
a reader of the project (Claude Opus~5.5 at effort high) checked the steps of
Theorem~\ref{thm:B} by hand and its identities with its own program. A
program by the writer (24 assertions) checks the charts of $S$ and $S_0$,
their derivatives, \eqref{eq:I} and $\sigma$; it and the run's certificate
are in the folder \texttt{checks} of the paper's workbench. On 25 November an
independent check (Claude Opus~5.5 at effort max) made a targeted check of
Lemma~\ref{lem:ext} and of Corollary~\ref{cor:C}, and found that $T$ is
singular at $o'$; its record is CHECK\_G7\_INDEPENDENT.md in the paper's
workbench. Before posting, the main session of the project (Claude Opus~5.5 at effort
max) read every page and checked Lemmas 2.1--2.6, the proof of Theorem~B and
Corollary~C line by line, including Lemma~2.6 and the identification of the
radial valuation in Corollary~C\,(ii).

\paragraph{The search of the literature.} On 22 and 25 November 2026 the
writer searched arXiv, the journal literature and the Hexagon catalogue for
wild covers, \'etale off the point, of surface singularities with a
non-rational exceptional curve; on 23 November the main session of this
project (Claude Opus~5.5 at effort max) read \cite{CS} at the source. None
of these contains the example. On 25 November the independent check read the
arXiv full texts of \cite{LS20,LS23,IS,Ma}: none has a wild cover over, or a
quotient singularity with, an exceptional curve of positive genus.
\cite{Ar75} (the publisher's site refused automated access) and \cite{Ar77}
(no free copy) were not opened. The covers of \cite{Ar75} have a regular
source, as recalled in \cite[Introduction]{LS20}, which forces rational
exceptional curves by Lemma~\ref{lem:rational}; the bases in \cite{Ar77} are
rational double points, whose exceptional curves are rational. So neither
contains the example. The full record is in PROVENANCE.md, in the
paper's workbench.

\begin{thebibliography}{SA12}
\bibitem[Ar75]{Ar75} M. Artin, \emph{Wildly ramified $\mathbb Z/2$ actions
in dimension two}, Proc. Amer. Math. Soc. 52 (1975), 60--64.
\bibitem[Ar77]{Ar77} M. Artin, \emph{Coverings of the rational double points
in characteristic $p$}, in: Complex Analysis and Algebraic Geometry,
Iwanami Shoten and Cambridge University Press, 1977, 11--22.
\bibitem[CS]{CS} S. D. Cutkosky, H. Srinivasan, \emph{Local fundamental
groups of surface singularities in characteristic $p$}, Comment. Math. Helv.
68 (1993), 319--332.
\bibitem[CST]{CST} J. Carvajal-Rojas, K. Schwede, K. Tucker,
\emph{Fundamental groups of $F$-regular singularities via $F$-signature},
Ann. Sci. \'Ec. Norm. Sup\'er. (4) 51 (2018), 993--1016; arXiv:1606.04088.
\bibitem[Ha]{Ha} R. Hartshorne, \emph{Complete intersections and
connectedness}, Amer. J. Math. 84 (1962), 497--508.
\bibitem[IS]{IS} H. Ito, S. Schr\"oer, \emph{Wild quotient surface
singularities whose dual graphs are not star-shaped}, arXiv:1209.3605.
\bibitem[LS20]{LS20} D. Lorenzini, S. Schr\"oer, \emph{Moderately ramified
actions in positive characteristic}, Math. Z. 295 (2020); arXiv:1904.08371.
\bibitem[LS23]{LS23} D. Lorenzini, S. Schr\"oer, \emph{Discriminant groups of
wild cyclic quotient singularities}, Algebra Number Theory 17 (2023);
arXiv:2101.08090.
\bibitem[Li78]{Li78} J. Lipman, \emph{Desingularization of two-dimensional
schemes}, Ann. of Math. (2) 107 (1978), 151--207.
\bibitem[Ma]{Ma} Y. Matsumoto, \emph{$\mu_p$- and $\alpha_p$-actions on K3
surfaces in characteristic $p$}, J. Algebraic Geom. 32 (2023), 271--322;
arXiv:1812.03466.
\bibitem[Mat]{Mat} H. Matsumura, \emph{Commutative ring theory}, Cambridge
University Press, 1986.
\bibitem[SA12]{SA12} GPT-6 Astra (OpenAI, ChatGPT Pro mode), run
RESEARCH\_SA12 of this project, final report of 20 November 2026 (603 pages),
Part~D.
\bibitem[Se]{Se} J.-P. Serre, \emph{Local fields}, Graduate Texts in
Mathematics 67, Springer, 1979.
\bibitem[Sti]{Sti} H. Stichtenoth, \emph{Algebraic function fields and
codes}, 2nd ed., Graduate Texts in Mathematics 254, Springer, 2009.
\end{thebibliography}

\end{document}
